\documentclass{report}
\usepackage{amsmath,amssymb}
\setlength{\parindent}{0mm}
\setlength{\parskip}{1em}
\begin{document}
\begin{center}
\rule{6in}{1pt} \
{\large Kenneth Esler \\
{\bf Fully Accelerating Classical AMG with GAMPACK (GPU Algebraic Multigrid PACKage)}}

Stone Ridge Technology \\ 2015 Emmorton Rd Suite 204 \\ Bel Air \\ MD 21015
\\
{\tt kesler@stoneridgetechnology.com}\\
Vincent Natoli\\
Yongpeng Zhang\\
Aleksander Samardzic\end{center}

Algebraic multigrid (AMG) has become the defacto standard for the
solution and preconditioning of difficult linear systems arising from
elliptic and parabolic PDEs in a broad range of domains. Its robust
convergence rates, linear complexity in system size, and ability to
handle unstructured and heterogeneous problems make it a nearly
indispensible tool in fields such as reservoir simulation.

In the last few years, the underlying computational architectures on
which high performance computing is performed have begun to undergo a
dramatic shift. As CPU clock frequencies have plateaued due to power
constraints, researchers have been forced to seek performance growth
through increasing parallelism. Multicore processors have provided an
evolutionary approach for which the methodologies based on coarse-grained
parallelism that were developed for distributed memory computers have
proved mostly adequate.

GPUs, however, require the simultaneous execution of thousands of
threads and a careful orchestration of memory access to achieve good
performance. This degree of multithreading, in turn, requires a more
fine-grained approach to parallelism. The regular access patterns and
naturally parallel smoothing methods of geometric multigrid has allowed
this method to benefit relatively easily from the increased memory
bandwidth and FLOP rate of GPUs. In contrast, the "branchy" execution and
largely random memory access patterns of AMG algorithms have made
acceleration with GPUs extremely challenging.

Initially, sparse matrix methods as a whole were largely dismissed by the
HPC community as inappropriate for acceleration on GPUs. Soon, however,
researchers proved this belief false by demonstrating significant
speedups by applying GPUs to iterative solvers with trivial or no
preconditioning. More complex preconditioners were gradually ported,
including those based on incomplete factorizations. Early work on
accelerating AMG focused solely on the solution stage, and again
dismissed AMG setup as impossible to accelerate. Later work demonstrated
that by selecting appropriate algorithms, constructing an AMG hierarchy
on GPU was possible, though the speedups achieved were more modest than
those typical of simpler algorithms.

In this talk, we discuss our efforts to fully accelerate AMG
algorithms on GPUs. By selecting appropriate algorithms and through
careful implementation, we find that it is possible to achieve a good
acceleration of classical AMG algorithms in both the setup and solve
stages. In particular, we find that partitioning the unknowns with a PMIS
approach and making use of distance-2 interpolators can provide a robust
hierarchy with ample fine-grained parallelism. Combining these algorithms
with an appropriate selection of smoothers and Krylov acceleration yields
a complete solver package with good speedup relative to available CPU AMG
solvers. We give examples, primarily drawn from reservoir simulation, and
compare performance with open-source CPU solutions. Finally, we discuss
our work to provide a solution which scales across multiple GPUs.


\end{document}
